% 统计图数据来自 2014 山东文科卷第 8 题题干及原始 PDF：
% 五个组距均为 1，频率/组距依次为 0.24, 0.16, 0.36, 0.16, 0.08。
% solid bars: [12,13), [13,14), [14,15), [15,16), [16,17].
% dashed guides: y=0.08, 0.16, 0.24, 0.36; x-axis is broken between 0 and 12.
\begin{tikzpicture}
    \begin{axis}[exam statistical histogram,
        width=8.4cm,
        height=5.4cm,
        xmin=0,
        xmax=7.35,
        ymin=0,
        ymax=0.42,
        axis lines=left,
        axis line style={-Stealth, line width=0.45pt},
        xtick=\empty,
        clip=false,
        ytick={0.08,0.16,0.24,0.36},
        yticklabels={0.08,0.16,0.24,0.36},
        ytick style={draw=none},
        yticklabel style={font=\normalsize, anchor=east, xshift=-2pt},
        x axis line style={-Stealth},
        y axis line style={-Stealth},
    ]
    \examstatxlabel{舒张压/kPa};
        % Source horizontal guides terminate at the matching bar boundary.
        \draw[dashed, line width=0.35pt] (axis cs:0,0.08) -- (axis cs:5,0.08);
        \draw[dashed, line width=0.35pt] (axis cs:0,0.16) -- (axis cs:3,0.16);
        \draw[dashed, line width=0.35pt] (axis cs:0,0.24) -- (axis cs:1,0.24);
        \draw[dashed, line width=0.35pt] (axis cs:0,0.36) -- (axis cs:3,0.36);

        % Broken x-axis leaves the long interval from 0 to 12 implicit.
        \examstatxbreak[black, line width=0.6pt]{0.43000000000000005}{0};

        % Histogram bars use the five frequencies read from the source chart.
        \addplot[/tikz/ybar interval,mark=none,black, line width=0.45pt, fill=white] coordinates {
          (1,0.24)
          (2,0.16)
          (3,0.36)
          (4,0.16)
          (5,0.08)
          (6,0)
        };

        \node[font=\normalsize, anchor=north,yshift=-4pt] at (axis cs:0,0.000000) {0};
        \node[font=\normalsize, anchor=north,yshift=-4pt] at (axis cs:1,0.000000) {12};
        \node[font=\normalsize, anchor=north,yshift=-4pt] at (axis cs:2,0.000000) {13};
        \node[font=\normalsize, anchor=north,yshift=-4pt] at (axis cs:3,0.000000) {14};
        \node[font=\normalsize, anchor=north,yshift=-4pt] at (axis cs:4,0.000000) {15};
        \node[font=\normalsize, anchor=north,yshift=-4pt] at (axis cs:5,0.000000) {16};
        \node[font=\normalsize, anchor=north,yshift=-4pt] at (axis cs:6,0.000000) {17};
        \examstatylabel{\(\dfrac{\text{频率}}{\text{组距}}\)};
    \end{axis}
\end{tikzpicture}
